% Noether P06 Korean producer draft — UNCHECKED
% T14-U191; ED0002 lines5627--5647; source 1969 LF B / EA89EE7E64AD0411AB1B8F9B072A4AA529DF8F1A462FD159FA241255D3510341
% Pointer v009 / ED0002: B06BE3530D9CF2E82B56FDBA7FE41D5D044DF2425DFA2A059D4939EAA2F7A6C2 / C9A125167ACB33D914EE4374B65AE7CDF0052F568371B8B77B720EA178ABF0E3
% Translation only; no review/build/render/approval. Lagrange genus-domain, elementary symmetry, resultant ±1, involution/Galois form, product index, and integral-rational combination remain held.

1. 제1종 상대정역 $\Gdom_{nn}$의 첫 번째 예로, \emph{주어진 순열군 $G$를 허용하는 $x_1\rcdots x_n$의 다항식 전체}, 곧 한 라그랑주 종영역의 다항식 전체를 고찰한다. $x_1\rcdots x_n$은 항등식
\[
\begin{array}{@{}l@{\quad}c@{}}
(1)& x_k^n-\sigma_1(x)x_k^{n-1}+\sigma_2(x)x_k^{n-2}+\cdots\pm\sigma_n(x)=0\quad(k=1,2,\ldots,n)
\end{array}
\]
에 따라 $\Gdom_{nn}$에 속하는 기본대칭함수들에 대하여 대수적으로 정수적이므로, 정의 V에 따라 $\Gdom_{nn}$에 대하여 대수적으로 정수적으로 종속된다. 따라서 $\Gdom_{nn}$은 제1종이다. 더 나아가 $\Gdom_{nn}$의 원소들이 동차형식이고 $n=\rho$이므로, $\Gdom_{nn}$은 제1종 영역으로서(\S\ 10의 결론에 따라) 한 정칙계를 이룬다. 실제로 (1)에 의해 기본대칭함수들의 결과식 $R_0$은 항등적으로 소멸할 수 없으며, 결과식 이론의 계산 규칙에 따르면 $R_0=\pm1$이다. $\Gdom_{nn}$의 \emph{인볼루션형식}은 --- $\Gdom_{nn}$이 제1종이고 $n=\rho$이므로 --- 대응하는 라그랑주 종영역의 인볼루션형식으로서 유일하게 정해진다. 이 영역에서 $x_1\rcdots x_n$에 켤레인 양들은 군 $G$의 순열들로 생기므로, 이 인볼루션형식은 다음과 같다.
\[
\begin{array}{@{}l@{\quad}c@{}}
(2)&
\begin{aligned}
  \Psi(z,u)&=\prod_i
  (z-u_1x_{i_1}-u_2x_{i_2}-\cdots-u_nx_{i_n})\\
  &=z^i+\sum{}'G_{\alpha\alpha_1\ldots\alpha_n}(x)
  z^\alpha u_1^{\alpha_1}\cdots u_n^{\alpha_n},
\end{aligned}
\end{array}
\]
여기서 곱은 $G$의 모든 순열에 걸쳐 취한다. (2)는 \glqq{}군 $G$의 갈루아형식\grqq{}을 나타낸다. 따라서 다음 사실이 성립한다.

\emph{순열군 $G$를 허용하는 $x_1\rcdots x_n$의 모든 다항식은 그 군의 갈루아형식의 계수 $G_{\alpha\alpha_1\ldots\alpha_n}(x)$들의 정수 유리 결합이다.}
